% Template LaTeX para TCC de graduação da UNIPAMPA — V2
% Copyright (c) 2026 Charles Quevedo Carpes.
% Autoria das contribuições próprias, organização e adaptações do template.
% This work may be distributed and/or modified under the conditions of the
% LaTeX Project Public License, version 1.3c or (at your option) any later version.
% License: https://www.latex-project.org/lppl/lppl-1-3c.txt (see LICENSE).
% This work has the LPPL maintenance status `author-maintained'.
% The Current Maintainer of this work is Charles Quevedo Carpes.
% Consulte AUTORIA.md para escopo, créditos e relação dos arquivos.

% Elemento opcional. Comente todo este arquivo em main.tex quando não houver glossário.
\chapter*{\raggedright GLOSSÁRIO}
\addtocontents{toc}{\protect\cftsetindents{chapter}{\cftlastnumwidth}{2em}}
\addcontentsline{toc}{chapter}{GLOSSÁRIO}

\begin{description}[style=multiline,leftmargin=4cm,labelwidth=3.5cm,font=\normalfont,itemsep=0.5\baselineskip]
  \item[Compilador] Programa responsável por transformar o código-fonte em um documento final.
  \item[LaTeX] Sistema de preparação de documentos baseado em marcação.
  \item[Template] Estrutura previamente configurada para padronizar um documento.
\end{description}
\clearpage